% Generated by maint/texbuild/theory_tex.py from theory/thought_experiments/demon_utm_four_noise_vectors.md.
% Edits here are overwritten. Edit the markdown, then run the script again.
\chapter{\texorpdfstring{Laplace's demon, a non-halting UTM, and the four noise vector magnitudes}{Laplace's demon, a non-halting UTM, and the four noise vector magnitudes}}
\label{chap:demon-utm-four-noise-vectors}

Combining Laplace's Demon (omniscient physical state knowledge) with a non-halting, mathematically perfect UTM allows us to treat video noise not as a random statistical distribution, but as a fully deterministic, computable function of atomic interactions, photon arrivals, and circuit thermodynamics.

Below is a framework representing this noise function by integrating across \(n\) still frames using four distinct vector magnitudes corresponding to fundamental physical noise components.

\section{\texorpdfstring{1. Theoretical Foundations (The Demon \& The UTM)}{1. Theoretical Foundations (The Demon \& The UTM)}}

Laplace's Demon provides the ground-truth micro-state \(S(x,y,z,t)\) for every atom in the sensor matrix, incident photons, and intervening medium.

The UTM acts as the universal evaluator, running the exact quantum-electrodynamic (QED) simulation to compute the ideal noise-free photon flux \(I(x, y, t)\) versus the observed corrupted pixel matrix \(F(x, y, t)\).

The residual noise tensor \(\mathcal{N}(x, y, t)\) is defined as:

\[
\mathcal{N}(x, y, t) = F(x, y, t) - I(x, y, t)
\]

\section{\texorpdfstring{2. The Four Noise Vector Magnitudes}{2. The Four Noise Vector Magnitudes}}

To integrate noise across a temporal sample of \(n\) still frames, we decompose \(\mathcal{N}\) into four orthogonal vector fields representing distinct physical noise sources, each quantified by its magnitude across the frame sequence:

\begin{enumerate}
\item \textbf{Photon Shot Noise Vector Magnitude} (\(\|\vec{v}_{\text{sh}}\|\)): captures the Poisson-distributed variance of discrete photon arrivals governed by quantum mechanics.
\item \textbf{Thermal/Read Noise Vector Magnitude} (\(\|\vec{v}_{\text{th}}\|\)): captures electron agitation and circuit-level voltage fluctuations within the sensor wells.
\item \textbf{Fixed-Pattern Noise Vector Magnitude} (\(\|\vec{v}_{\text{fp}}\|\)): captures spatial manufacturing inhomogeneities and dark current variations across individual sensor pixels.
\item \textbf{Quantization/Stochastic Vector Magnitude} (\(\|\vec{v}_{\text{q}}\|\)): captures analog-to-digital conversion rounding errors and high-frequency atmospheric scintillation.
\end{enumerate}

\section{\texorpdfstring{3. The Integrated Noise Function}{3. The Integrated Noise Function}}

Let each frame at time index \(t \in \{1, 2, \dots, n\}\) be represented over spatial coordinates \((x, y)\). We integrate the noise function \(\Psi\) by summing the inner products of these four vector magnitudes weighted across the \(n\)-frame temporal window:

\[
\Psi(x, y, n) = \sum_{t=1}^{n} \left( \|\vec{v}_{\text{sh}}(x,y,t)\| + \|\vec{v}_{\text{th}}(x,y,t)\| + \|\vec{v}_{\text{fp}}(x,y)\| + \|\vec{v}_{\text{q}}(x,y,t)\| \right) \cdot w(t)
\]

where \(w(t)\) is a temporal weighting function computable by the UTM to account for shutter exposure profiles and sensor thermal accumulation over the sample duration \(n\).

Because Laplace's Demon tracks the exact phase space of every particle, the “random” variables driving these four vectors are reduced to deterministic initial conditions, allowing the UTM to evaluate \(\Psi(x, y, n)\) with zero statistical variance.
